\origpage{148}{ROBERTSON AND SEYMOUR}
\hspace*{12pt}(6) $d(z, c) \le \lambda + 1$.

\textit{Subproof.} Let $f \in E(H')$ be such that $\eta'(f)$ is the edge of $Q$ incident with $\eta(c)$. If $f \not\subseteq \Lambda$, then $f \in E(H)$ and $\eta(f) = \eta'(f) \in E(Q)$, a contradiction since no edge of $Q$ is in $E(\eta(H))$. Thus $f \subseteq \Lambda$, and so $c \in \bar{\Lambda}$. Since $c \in V(H)$ there exists $\sigma \in \Lambda$ with $d(c, \sigma) \le 1$, and consequently $d(z, c) \le d(z, \sigma) + 1 \le \lambda + 1$. This proves (6).

From (6) we deduce that $c \ne a$, since $d(z, a) > \lambda' \ge \lambda + 1$, and so $P \cup Q$ is an $\eta(H)$-path with ends $\eta(a), \eta(c)$. From (2) and (6),
\begin{equation*}
d(a, c) \ge d(z, a) - d(z, c) \ge (\lambda + \mu + 3) - (\lambda + 1) = \mu + 2 \ge \mu,
\end{equation*}
and so $H, \eta, \mathscr{T}$ is $\mu$-stepped and not $(\lambda + \mu + 2, \mu)$-level. This completes the proof. $\square$

From (8.2) and (8.3) we obtain

\stmt
(8.4) \textit{For any surface $\Sigma$ with $bd(\Sigma) = \emptyset$, and any integers $p, \phi, \mu \ge 0$, there are integers $\kappa, \theta, \lambda, \rho \ge 0$ such that the following holds. Let $\mathscr{T}^*$ be a tangle in a graph $G$ such that there is no $\mathscr{T}^*$-central segregation of $G$ of type $(\rho, \kappa)$ with an arrangement in $\Sigma$. Then either}

\hspace*{12pt}(i) \textit{$\mathscr{T}^*$ controls a $K_p$ minor of $G$, or}

\hspace*{12pt}(ii) \textit{there is a $\Sigma'$-span of order $\ge \phi$, where $\Sigma'$ is a surface obtained by adding a crosscap to $\Sigma$, or}

\hspace*{12pt}(iii) \textit{no $\Sigma$-span of order $\ge \theta$ is $(\lambda, \mu)$-level.}

\textit{Proof.} Choose $\kappa, \theta, \lambda', \rho$ to satisfy (8.2). Let $\lambda = \lambda' + \mu + 2$; and we claim that (8.4) holds. For let $\mathscr{T}^*, G$ be as (8.4). By (8.2), we may assume that no $\Sigma$-span of order $\ge \theta$ is $(\lambda', \mu)$-flat. By (8.3), no $\Sigma$-span of order $\ge \theta$ is $(\lambda, \mu)$-level, as required. $\square$

\sectionhead{\textsc{References}}

1. N. Robertson and P. D. Seymour, Graph minors. IX. Disjoint crossed paths, \textit{J. Combin. Theory Ser. B} \textbf{49} (1990), 40--77. 2. N. Robertson and P. D. Seymour, Graph minors. X. Obstructions to tree-decomposition, \textit{J. Combin. Theory Ser. B} \textbf{52} (1991), 153--190. 3. N. Robertson and P. D. Seymour, Graph minors. XI. Circuits on a surface, \textit{J. Combin. Theory Ser. B} \textbf{60} (1994), 72--106. 4. N. Robertson and P. D. Seymour, Graph minors. XII. Distance on a surface, \textit{J. Combin. Theory Ser. B} \textbf{64} (1995), 240--272. 5. N. Robertson and P. D. Seymour, Graph minors. XIV. Extending an embedding, \textit{J. Combin. Theory Ser. B} \textbf{65} (1995), 23--50.
